### Conventions / definitions

* All graphs are **finite, simple, undirected**: (G=(V,E)) with no loops or multiple edges.
* (|V|=n).
* (K_r)-free means (G) contains **no** complete subgraph on (r) vertices.
* The **average degree** is
  [
  t:=\frac{1}{n}\sum_{v\in V}d(v)=\frac{2|E|}{n}.
  ]
* An **independent set** is a set (I\subseteq V) with no edges inside (I); its maximum size is (\alpha(G)).
* (\log) is the **natural logarithm**; changing the base only changes constants, which are absorbed into (\gg_r).
* The notation (A\gg_r B) means: there exists a constant (c_r>0) depending only on (r) such that (A\ge c_r B) for all admissible parameters.

### Literal statement (as written)

> For fixed (r), every (K_r)-free (n)-vertex graph of average degree (t) has an independent set of size (\gg_r \frac{\log t}{t}n).

**Ambiguity/misstatement:** if (t=0) (edgeless graph), then (\log t) is undefined. Since average degree can be (0), the statement is not literally well-formed without an additional convention.

### Minimal corrected statement (standard)

Fix an integer (r\ge 3).
There exists a constant (c_r>0) such that for every (n\in\mathbb N) and every (K_r)-free graph (G) on (n) vertices with average degree (t\ge 2),
[
\alpha(G);\ge; c_r,\frac{\log t}{t},n.
]
(For (t<2) the bound is anyway trivial up to constants, because (\alpha(G)\ge n/(t+1)\ge n/3) while ((\log t)/t) is (O(1)).)

This is the precise version I will analyze.

---
